\section{Discussion and limitations}
\label{sec:discussion}

\paragraph{What the evidence says.} The capability-saturation principle makes two claims: coordination helps while the SAS is weak, and its value declines as the SAS grows stronger. On the two cells that test the first claim in this setting, the MAS loses by 13 and 22 points; we find no support for the prediction in this setting, for this architecture. The one gain in the study appears on the straddling cell under a second simulator. The second claim is left undecided: the slope test detects no dependence, and its interval is too wide to exclude a steep decline, because model-by-domain effects dominate the variation between cells. The mechanism that motivates a MAS here, catching wrongful writes, shows no detectable effect. For systems that serve users under a written policy, the SAS was the better choice at every capability level we measured under a fixed simulator, and at a fraction of the cost. These results are for the MAS as built. Its planner never saw the tool schemas, and \S\ref{sec:compliance} reports how far this moved its orders away from what the database agent executed. A MAS design for such settings therefore needs compute-matched evidence from the setting itself, since results from autonomous benchmarks do not carry over.

\paragraph{Why conversation with a user penalises coordination.} The user holds information the agent needs, on telecom part of the control as well, and the benchmark is half-duplex, so every tool call passes through the user-facing agent's turn and the orchestrator plans before each assignment. The user supplies information turn by turn, so the task cannot be decomposed in advance, and decomposable tasks are where coordination drew its largest benefits in the original study \citep{kim2026capable}. The reward is a state at the end of a long conversation, so a MAS that spends its ceiling early fails outright.

\paragraph{Limitations.} (1) One architecture is tested, and it is reported as built. Its orchestrator plans without the tool schemas that only the database agent holds, so under the stricter audit of \S\ref{sec:compliance} only 25\% of its write orders were executed as issued, against the 90\% the plan required; the pre-registered fidelity bar fails in every cell and the result is a labelled deviation. A fix that shows the orchestrator the schemas raised the as-issued share from 25\% to 72\% in a three-episode smoke test, but the grid was not re-run with it, so the paper measures the architecture with its planner blind to the schemas. A null below 45\% is stronger evidence against transfer than the original best-of-four validation gives, and a null above 45\% is weaker. (2) Only two cells lie clearly below 45\%, so the principle's distinctive prediction rests on two cells and the slope is identified mostly above the threshold. (3) The retail \texttt{qwen3.5-4b} cell was run on a later model build than the one registered, so it sits outside the fits; it shows a 51-point MAS loss at 74\% SAS accuracy (Appendix~\ref{app:local4b}). The airline \texttt{qwen3.5-4b} MAS was run by a collaborator on the registered build against a budget taken from an S1 on a different software stack; paired with the September S1 it is reported as a sensitivity whose sign depends on how that S1's crashed episodes are scored (Appendix~\ref{app:local4b}). (4) The user is a model, and simulators are more cooperative than people, so where a cell sits relative to 45\% may differ with human customers; swapping the simulator already turns the straddling cell's tie into a 10-point MAS gain (R2). (5) The official reward is blind to verbal violations on airline and telecom; we measure this blind spot on airline only, commission flags are classified by a fixed text rule, and the judge was not checked against human labels. (6) The models postdate the benchmark's release and may have seen it in training, which inflates SAS and MAS scores alike; the 4-billion open-weight model's telecom scores (94.7\% and 95.3\% across builds) sit above every closed model's, which is consistent with contamination though we cannot test it. (7) On the local route the benchmark passes each trial's seed to the model, so one seed value fixed a long-reasoning behaviour in the MAS orchestrator across three local cells; the effect is confirmed by a replay and the cells are reported per trial, with no trial excluded (Appendix~\ref{app:local4b}).


\paragraph{Further limitations.} (8) The OpenAI snapshots are scheduled for retirement on 11 December 2026; after that date the released trajectories are the record. (9) The 56\% plan-parse failure rate of \texttt{qwen3.5-4b} is a property of that model under our parser; a parser that completes a missing closing brace is an obvious post hoc variant and was not run. (10) A write followed by a compensating revert leaves the database unchanged and is invisible to the primary measure; three such candidates in the pilot run are listed for hand adjudication. (11) The reasoning-only sensitivity meter cannot split 1.2\% of gate-run output tokens, whose calls carry no details block; the primary meter still counts them. (12) The judge model is pinned; the secondary measure depends on it. (13) The uncapped arm uses the three cells at or below 45\% and the extension the same three, each with three seeds; in the uncapped arm, excluding the 6 telecom $\times$ \texttt{gpt-5-nano} episodes in which plan text reached the customer moves that cell's gain from $-10.0$ to $-12.0$. (14) The \texttt{qwen3.5-9b} \Sone\ values were computed by a collaborator; the per-episode files needed for the paired intervals of two of its cells are not held by us. (15) The open-weight rungs run on non-standard stacks: a local 4-bit server, and an unpinned hosted route that did not enforce the cap on reasoning tokens. (16) The hosted APIs drop the per-trial seed, so the closed-model cells carry no seed effect of the kind found on the local route; whether the SAS on the local route carries a seed effect of its own was not tested. (17) The fix for the fidelity bar was smoke-tested on 5 tasks in each of three cells; its effect on the pass rate was not measured, and the two misses that were policy refusals by the database agent show that no orchestrator-side change can reach the 90\% bar on its own. (18) The collaborator's September SAS runs of \texttt{qwen3.5-4b} on airline and retail, on a different serving stack, scored 39.7 and 36.3\% with 45 to 48\% of episodes ending in agent errors that our stack did not show, so whether the airline cell would sit below 45\% on our stack is unknown.
